\chapter{System Study}

Evaluating modern sports technology requires analyzing existing analytical workflows and identifying systemic deficiencies. The system study phase establishes the rationale for CricketVision by examining legacy cricket tracking setups, reviewing academic research in sports science and computer vision, conducting a gap analysis, and specifying software and hardware requirements.

\section{Existing System}

\subsection{Overview and Limitations}
Existing cricket management setups typically rely on commercial sports platforms (e.g., Cricinfo, Cricbuzz, standalone media player software, and offline Excel spreadsheets). While these tools provide baseline scoring metrics, they exhibit severe operational constraints when applied to elite franchise management:

\begin{enumerate}
    \item \textbf{Static Historical Data:} Existing scorecards display aggregate runs and wickets but fail to segregate performance across distinct match phases (Powerplay overs 1--6, Middle overs 7--15, Death overs 16--20).
    \item \textbf{No Integrated Biomechanics:} Video review in standard software lacks computer-vision pose tracking, requiring expensive offline motion-capture hardware (e.g., Vicon, Qualisys optical systems costing hundreds of thousands of dollars).
    \item \textbf{Absence of Real-Time Win Probability Simulation:} Conventional systems cannot simulate match outcomes in real time when key variables change (e.g., sudden loss of wickets, dew factor shifts, weather adjustments).
    \item \textbf{Coarse Fatigue Management:} Workload tracking is usually recorded manually in physical logs or generic GPS vests, making it difficult to prevent lumbar stress fractures or shoulder impingements.
    \item \textbf{No Persona-Based Access Control:} Standard tools lack tailored permissions; coaches, players, and analysts view identical generic interfaces without context-aware personalization.
\end{enumerate}

\begin{table}[H]
\centering
\caption{Comparison of Existing Systems vs.\ CricketVision Platform}
\label{tab:system_comparison}
\begin{tabular}{p{3.2cm} p{4.5cm} p{5cm}}
\toprule
\textbf{Feature / Metric} & \textbf{Conventional Tools} & \textbf{CricketVision Platform} \\
\midrule
Data Storage & Static Flat Files / CSVs & Centralized MongoDB (100+ Players) \\
Phase-Wise Stats & Manual Calculation Only & Automated (Powerplay, Middle, Death) \\
Match Simulation & Static DLS Lookups & Dynamic Ball-by-Ball Win Prob Engine \\
Visual Analytics & Basic Bar/Pie Charts & Recharts Radar, 360\textdegree{} Wagon Wheel SVG \\
Biomechanics & Manual Video Frame Pause & Computer Vision Pose Angle Extraction \\
Role Access & Uniform Single View & Role-Based Auth (Coach, Player, Analyst) \\
AI Tactical Advice & Static PDF Scouting Reports & Interactive Natural Language AI Coach \\
\bottomrule
\end{tabular}
\end{table}

\subsection{Literature Survey}
Recent academic literature in sports analytics, computer vision, and predictive modeling emphasizes the transition toward automated, data-driven frameworks:

\begin{enumerate}
    \item \textbf{Perera et al.\ (2018) --- \textit{Sports Analytics in Cricket: A Survey}}:\\
    Highlights that traditional aggregate metrics (e.g., batting average, bowling economy) fail to capture context-specific pressure situations. The authors advocate for phase-based efficiency ratings, venue-adjusted performance metrics, and clutch situational indices. Their research directly motivates CricketVision's phase-wise breakdown modules.

    \item \textbf{Moorthy \& DeSarbo (2020) --- \textit{Predictive Modeling in Franchise Cricket}}:\\
    Demonstrates that stochastic simulation models incorporating pitch condition, weather parameters, and bowler-batter matchups yield significantly higher predictive accuracy for match outcomes compared to static linear regression models. This work inspired CricketVision's logistic win-probability engine incorporating RRR decay, wicket survival, and dew factors.

    \item \textbf{Cao et al.\ (2021) --- \textit{Real-time Multi-Person 2D Pose Estimation using Part Affinity Fields}}:\\
    Establishes the foundation for markerless biomechanics tracking in sports, proving that joint location coordinates derived from video frames can accurately measure body segment angles (such as elbow flex and shoulder inclination) without physical motion sensors. This forms the theoretical basis of CricketVision's VideoAnalyzerView pose engine.

    \item \textbf{Ferdinands et al.\ (2019) --- \textit{Biomechanical Analysis of Lumbar Loading in Cricket Fast Bowlers}}:\\
    Explored the kinematic factors associated with bowling performance and lumbar spinal injury risk. The research demonstrated that front knee flexion at impact and elbow extension velocity are critical determinants of both ball speed and injury susceptibility, reinforcing the need for continuous automated video monitoring.
\end{enumerate}

\section{Proposed System}

\subsection{Key Highlights and Capabilities}
The proposed \textbf{CricketVision} platform introduces a full-stack solution designed to address the limitations of existing tools:

\begin{itemize}
    \item \textbf{MERN Architecture \& Database Seeding:} Built on Node.js and Express.js with Mongoose ODM connected to MongoDB. Automatically populates a dataset of 100+ international and IPL players on initial deployment.
    \item \textbf{Role-Based Authentication Engine:} Includes pre-configured persona presets (e.g., Coach Dravid, Player Kohli / Bumrah, Analyst Alex Morgan) enforcing access rights across sensitive endpoints.
    \item \textbf{Interactive 360\textdegree{} Wagon Wheel \& Pitch Heatmap:} Utilizes SVG polar-to-Cartesian math to calculate shot directions across 8 stadium sectors alongside pitch length distribution heatmaps.
    \item \textbf{Stochastic Match Simulator:} Executes ball-by-ball win-probability calculations considering target score, current run rate, required run rate, remaining wickets, weather impact factors, and pitch friction coefficients.
    \item \textbf{Pose Estimation Biomechanics Video Analyzer:} Features an interactive video engine rendering pose tracking skeletons, overlaying elbow flex angles and stride length metrics, and verifying ICC 15\textdegree{} compliance.
    \item \textbf{AI Coach Assistant:} Provides a contextual natural language strategy engine suggesting optimal bowling rotations, field placements, rehabilitation drills, and team selection tweaks.
\end{itemize}

\section{Requirement Analysis}

\subsection{S/W Requirement}
\begin{table}[H]
\centering
\caption{Software Requirements Specification}
\label{tab:software_reqs}
\begin{tabular}{p{3.5cm} p{4cm} p{4.5cm}}
\toprule
\textbf{Component Layer} & \textbf{Technology} & \textbf{Specification / Version} \\
\midrule
Operating System & Windows 10/11, Ubuntu 22.04, macOS & 64-bit OS \\
Runtime Environment & Node.js (V8 Engine) & Node.js v18.16.0+ \\
Web Framework & Express.js RESTful API & v4.18.2+ \\
Database Server & MongoDB Community + Mongoose ODM & MongoDB v6.0+ \\
Frontend Library & React 18 (Hooks, Context) & v18.2.0+ \\
Build Tooling & Vite + HMR & v4.4.0+ \\
UI Styling & Tailwind CSS + PostCSS & v3.3.0+ \\
Data Visualization & Recharts + Lucide Icons & Recharts v2.7.0+ \\
Language Support & JavaScript ES2022+, Python 3 & Python 3.10+ \\
Version Control & Git \& GitHub & Git v2.40+ \\
API Testing & Postman, Chrome DevTools & Postman v10.x \\
\bottomrule
\end{tabular}
\end{table}

\subsection{H/W Requirement}
\begin{table}[H]
\centering
\caption{Hardware Requirements Specification}
\label{tab:hardware_reqs}
\begin{tabular}{p{3.5cm} p{4cm} p{4.5cm}}
\toprule
\textbf{Hardware Component} & \textbf{Development Machine} & \textbf{Client / End-User} \\
\midrule
Processor (CPU) & Intel Core i5/i7, AMD Ryzen 5/7 & Dual-Core 2.0 GHz+ \\
System Memory (RAM) & 16 GB DDR4 / DDR5 & 8 GB RAM minimum \\
Storage Space & 512 GB NVMe SSD & 10 GB free disk space \\
Display Resolution & 1920$\times$1080 Full HD & 1366$\times$768 or higher \\
Network Interface & Gigabit Ethernet / Wi-Fi 6 & Broadband (10 Mbps+) \\
GPU (Optional) & Dedicated for video pose rendering & Integrated graphics sufficient \\
\bottomrule
\end{tabular}
\end{table}
